\documentclass[a4paper,12pt]{article}

\usepackage[T2A]{fontenc}
\usepackage[utf8]{inputenc}
\usepackage[russian]{babel}
\usepackage{amsmath,amssymb,amsthm,mathtools}
\usepackage{helvet}
\renewcommand{\familydefault}{\sfdefault}
\usepackage{icomma}
\usepackage{geometry}
\usepackage{microtype}
\usepackage{tikz}
\usetikzlibrary{calc,arrows.meta,positioning,shapes.geometric}
\usepackage{tabularx,booktabs}
\usepackage{enumitem}
\usepackage{hyperref}
\usepackage{parskip}
\usepackage{fancyhdr}
\usepackage{setspace}
\usepackage{xcolor}

\geometry{margin=2.2cm}
\onehalfspacing

\definecolor{accent}{HTML}{008080}
\definecolor{darkaccent}{HTML}{004D4D}
\definecolor{textgray}{HTML}{333333}

\hypersetup{
colorlinks=true,
linkcolor=darkaccent
}

\pagestyle{fancy}
\fancyhf{}
\fancyfoot[C]{\small Лёвин Артём Александрович эксклюзивно для Даниила Кичука}
\fancyfoot[R]{\thepage}

\setlength{\parindent}{0pt}

\begin{document}

\begin{titlepage}
\begin{center}

\vspace*{2cm}

{\Large\textcolor{darkaccent}{\textbf{Чек-лист}}}

\vspace{1cm}

{\Huge\textbf{Степени и показательные уравнения}}

\vspace{1.5cm}

{\Large Даниил Кичук}\\[0.3cm]
{\large ЕГЭ профильная математика}

\vfill

{\large Лёвин Артём Александрович}\\
{\large \today}

\vspace{1cm}

\begin{tikzpicture}
\draw[accent,line width=1pt]
(0,0) .. controls (2,0.8) and (5,-0.4) .. (9,0.5);
\end{tikzpicture}

\end{center}
\end{titlepage}

\section*{1. Главная идея темы}

В задачах со степенями важно начинать не с формулы, а с поиска правильного направления.

\textbf{Концепция темы:}

\[
\boxed{\text{привести выражение к степенному виду}}
\]

После этого становятся доступны свойства степеней.

В первую очередь обращаем внимание на:

\begin{itemize}
\item составные числа;
\item корни;
\item дроби.
\end{itemize}

Например:

\[
8=2^3,\qquad
\sqrt{8}=2^{3/2},\qquad
\frac1{32}=2^{-5}.
\]

\section*{2. Свойства степеней}

\[
a^m\cdot a^n=a^{m+n}
\]

Условие:
\[
\text{основания должны совпадать}.
\]

\[
\frac{a^m}{a^n}=a^{m-n},
\qquad a\neq0
\]

\[
(a^m)^n=a^{mn}
\]

\[
(ab)^n=a^nb^n
\]

\[
\left(\frac ab\right)^n=\frac{a^n}{b^n},
\qquad b\neq0
\]

\[
a^{-n}=\frac1{a^n},
\qquad a\neq0
\]

\[
\sqrt[n]{a^m}=a^{\frac mn}
\]

Главный ориентир:

\[
\boxed{
\text{внутренняя степень делится на внешнюю}
}
\]

Например:

\[
\sqrt[3]{2^5}=2^{5/3}.
\]

\section*{3. Разложение составных чисел}

\textbf{Простое число} имеет ровно два делителя:
\[
1 \text{ и само число}.
\]

Например:

\[
7 \text{ --- простое},\qquad
9 \text{ --- составное}.
\]

Число $1$ не является ни простым, ни составным.

Чтобы работать со степенями, составные числа раскладываем на простые множители.

Пример:

\[
72=2\cdot36=2\cdot2\cdot18
=2\cdot2\cdot2\cdot9
=2^3\cdot3^2.
\]

Алгоритм:

\begin{enumerate}
\item найти самый маленький делитель;
\item разделить число;
\item повторять процесс до получения единицы.
\end{enumerate}

\section*{4. Перевод корней в степени}

Пример:

\[
\sqrt8
=
\sqrt{2^3}
=
(2^3)^{1/2}
=
2^{3/2}.
\]

Другой пример:

\[
\sqrt[4]{81}
=
\sqrt[4]{3^4}
=
3.
\]

\section*{5. Перевод дробей в степени}

Используем правило:

\[
\frac1{a^n}=a^{-n},
\qquad a\neq0.
\]

Пример:

\[
\frac1{64}
=
\frac1{2^6}
=
2^{-6}.
\]

\section*{6. Типовой подход к вычислениям}

При решении выражения:

\[
\frac{5^{0.36}}{25^{0.32}}
\]

действуем по алгоритму:

\begin{enumerate}
\item привести числа к одному основанию;
\item применить свойства степеней;
\item выполнить вычисления.
\end{enumerate}

Так как

\[
25=5^2,
\]

получаем:

\[
\frac{5^{0.36}}{(5^2)^{0.32}}
=
\frac{5^{0.36}}{5^{0.64}}
=
5^{-0.28}.
\]

\section*{7. Показательные уравнения}

Главная цель:

\[
a^{f(x)}=a^{g(x)}.
\]

Только если основания одинаковые, можно сравнить показатели:

\[
f(x)=g(x).
\]

Пример:

\[
2^{x+3}=8.
\]

Представим правую часть через основание $2$:

\[
8=2^3.
\]

Тогда:

\[
2^{x+3}=2^3.
\]

Следовательно:

\[
x+3=3,
\]

\[
x=0.
\]

\section*{8. Когда нельзя сразу сравнивать степени}

Нельзя делать переход:

\[
2^x=3^x
\Rightarrow x=x
\]

так как основания разные.

Сначала необходимо привести выражение к единому основанию.

Например:

\[
2^x\cdot8=2^{x+5}
\]

Так как

\[
8=2^3,
\]

получаем:

\[
2^x\cdot2^3=2^{x+5}.
\]

После преобразования:

\[
2^{x+3}=2^{x+5},
\]

что приводит к сравнению показателей.

\section*{9. Типичные ошибки}

\begin{itemize}
\item складывать показатели при разных основаниях;
\item забывать условия применения формул;
\item неверно работать с отрицательной степенью;
\item сокращать выражения без проверки нулей в знаменателе;
\item терять скобки при переносах.
\end{itemize}

Перед каждым преобразованием задавайте вопрос:

\[
\boxed{\text{какую формулу я сейчас могу применить и почему?}}
\]

\clearpage

\section*{Тренировочный блок}

\begin{enumerate}[itemsep=1em]

\item Представьте в степенном виде:
\[
\sqrt{32}
\]

\item Представьте в степенном виде:
\[
\frac1{128}
\]

\item Разложите на простые множители:
\[
108
\]

\item Представьте в степенном виде:
\[
\sqrt[3]{81}
\]

\item Упростите:
\[
\frac{3^7}{3^4}
\]

\item Вычислите:
\[
\frac{2^{3/2}}{2^{1/2}}
\]

\item Решите уравнение:
\[
5^{x+1}=125
\]

\item Решите уравнение:
\[
3^{2x-1}=27
\]

\item Упростите:
\[
\frac{4^{x}}{2^{x}}
\]

\item Решите уравнение:
\[
2^{x+2}=4\sqrt2
\]

\end{enumerate}

\clearpage

\section*{Ответы}

\renewcommand{\arraystretch}{1.3}

\begin{center}
\begin{tabularx}{\linewidth}{cX}
\toprule
№ & Ответ\\
\midrule
1 & $2^{5/2}$\\
2 & $2^{-7}$\\
3 & $108=2^2\cdot3^3$\\
4 & $3^{4/3}$\\
5 & $3^3=27$\\
6 & $2$\\
7 & $x=2$\\
8 & $x=2$\\
9 & $2^x$\\
10 & $x=\frac12$\\
\bottomrule
\end{tabularx}
\end{center}

\clearpage

\end{document}